\documentclass[11pt,a4paper]{article}

% -------------------- Langue / encodage --------------------
\usepackage[T1]{fontenc}
\usepackage[french]{babel}
\usepackage{array}

% -------------------- Mise en page --------------------
\usepackage{geometry}
\geometry{
  a4paper,
  landscape,
  left=1.6cm,
  right=1.6cm,
  top=1.0cm,
  bottom=1.2cm
}

% -------------------- Maths / mise en forme --------------------
\usepackage{amsmath,amssymb}
\usepackage{parskip}
\usepackage[explicit]{titlesec}
\usepackage[normalem]{ulem} % pour \uline sans casser \emph
\usepackage{titlesec}

\usepackage{fancyhdr}
\usepackage{lastpage}
\usepackage[hidelinks]{hyperref}
\usepackage{enumitem}
\usepackage[table]{xcolor} % <-- IMPORTANT : pour \rowcolor dans les tableaux
\usepackage{paracol}



% -------------------- TikZ --------------------
\usepackage{tikz}
\usetikzlibrary{arrows.meta,calc}

% -------------------- Variables --------------------
\newcommand{\FicheTitre}{Compléments de dérivation -- Fiche de cours}
\newcommand{\Niveau}{Terminale générale -- Mathématiques}
\newcommand{\Annee}{Année scolaire 2025/2026}

\newcommand{\SiteURL}{https://naoufel-coursparticuliers.netlify.app/}
\newcommand{\SiteTexte}{naoufel-coursparticuliers.netlify.app}

% -------------------- Footer --------------------
\pagestyle{fancy}
\fancyhf{}
\lfoot{\textit{\Niveau \;--\; \Annee}}
\cfoot{\thepage/\pageref{LastPage}}
\rfoot{\href{\SiteURL}{\texttt{\SiteTexte}}}
\renewcommand{\headrulewidth}{0pt}
\renewcommand{\footrulewidth}{0.4pt}

% -------------------- Spacing --------------------
\setlength{\parskip}{4pt}
\setlength{\parindent}{0pt}
\setlist[itemize]{leftmargin=*, itemsep=2pt, topsep=2pt}

% -------------------- Titres --------------------
%\titleformat{\section}{\large\bfseries}{\thesection.}{0.5em}{} %noir
%\titleformat{\subsection}{\normalsize\bfseries}{\alph{subsection}.}{0.5em}{}\titleformat{\section}
% --- TITRES colorés + petite ligne sous le TEXTE (sans souligner le numéro) ---
% SECTION : numéro rouge (non souligné) + titre rouge souligné
\titleformat{\section}[hang]
  {\normalsize\bfseries}
  {\textcolor{red}{\thesection.}}
  {0.2em}
  {\textcolor{red}{\uline{#1}}}

% SUBSECTION : lettre verte (non soulignée) + titre vert souligné
\titleformat{\subsection}[hang]
  {\normalsize\bfseries}
  {\hspace*{1.2em}\textcolor{green!60!black}{\thesubsection}}
  {0.4em}
  {\textcolor{green!60!black}{\uline{#1}}}


% -------------------- Bandeau titre --------------------
\newcommand{\TitreBandeau}[1]{%
  \noindent\hrule
  \vspace{0.55em}
  \begin{center}
    {\LARGE\bfseries #1}
  \end{center}
  \vspace{0.55em}
  \noindent\hrule
  \vspace{0.9em}
}

% -------------------- Helpers --------------------
\newcommand{\defi}{\textbf{Définition :} }
\newcommand{\prop}{\textbf{Propriété :} }
\newcommand{\rem}{\textbf{Remarque :} }

% ============================
% Graphique : Point d'inflexion (coloré + tangente cadrée + axes titrés)
% ============================
\newcommand{\GraphInflexion}{%
\begin{center}
\resizebox{0.78\linewidth}{!}{%
\begin{tikzpicture}[>=Latex, every node/.style={font=\small}]

% Fenêtre
\def\xmin{0} \def\xmax{5}
\def\ymin{0} \def\ymax{4}

% Fonds colorés (concave / convexe)
\fill[blue!15] (\xmin,\ymin) rectangle (3,\ymax);
\fill[pink!15] (3,\ymin) rectangle (\xmax,\ymax);

% Grille
\draw[step=0.5, very thin, gray!35] (\xmin,\ymin) grid (\xmax,\ymax);

% Axes
\draw[thick,->] (\xmin,0) -- (\xmax+0.2,0) node[below right] {$x$};
\draw[thick,->] (0,\ymin) -- (0,\ymax+0.2) node[above left] {$y$};

% Ligne verticale (abscisse du point)
\draw[dashed] (3,\ymin) -- (3,\ymax);

% Tout ce qui suit est "cadré" (tangente ne déborde pas)
\begin{scope}
\clip (\xmin,\ymin) rectangle (\xmax,\ymax);

% Courbe (rouge)
\draw[very thick, red, domain=\xmin:\xmax, samples=250]
plot (\x, {0.2*(\x-3)^3 - (\x-3) + 2.3});

% Tangente (violette) : y = - (x-3) + 2.3
\draw[thick, violet, domain=\xmin:\xmax]
plot (\x, {-1*(\x-3) + 2.3});

\end{scope}

% Point A
\fill[blue] (3,2.3) circle (3.2pt);
\node[blue, above right] at (3,2.1) {$A(a;f(a))$};

% Labels zones
\node at (1.2,2.2) {\normalsize Concave};
\node at (4.0,3.0) {\normalsize Convexe};

% Label Cf (positionné sur la courbe)
\node[red] at (4.15,1.8) {\Large $\mathcal{C}_f$};

% Label Ta 
\node[violet] at (2,3.6) {\Large $\mathcal{T}_a$};

\end{tikzpicture}
}%
\end{center}
}



\begin{document}

\TitreBandeau{\FicheTitre}

% -------------------- 2 colonnes --------------------
\columnratio{0.5,0.5}
\setlength{\columnsep}{1.2cm}
\setlength{\columnseprule}{0.6pt}
\def\columnseprulecolor{\color{black}}

\begin{paracol}{2}

% =====================================================
% ================= COLONNE GAUCHE ====================
% =====================================================

\section{Composition de fonctions}

\subsection{Définition}
\defi Soient $u$ une fonction définie sur $I$ et $f$ une fonction définie sur un intervalle $J$ contenant les valeurs de $u$.

On appelle \textbf{fonction composée} de $u$ par $f$ la fonction notée $f\circ u$ définie sur $I$ par :
\[
(f\circ u)(x)=f(u(x)).
\]

\subsection{Dérivée (formule de chaîne)}
\prop Si $u$ est dérivable sur $I$ et $f$ est dérivable sur $J$, alors $f\circ u$ est dérivable sur $I$ et :
\[
(f\circ u)'(x)=u'(x)\,f'(u(x)).
\]

\rem On dérive la fonction extérieure, puis on multiplie par la dérivée de la fonction intérieure.

\subsection{Formules usuelles avec une fonction $u$}

Si $u$ est dérivable :

\[
(e^u)' = u'e^u
\qquad
(\ln(u))'=\frac{u'}{u} \ (u>0)
\]

\[
(u^n)' = n u' u^{\,n-1}
\qquad
\left(\frac{1}{u}\right)' = -\frac{u'}{u^2}
\]

\[
(\sqrt{u})'=\frac{u'}{2\sqrt{u}} \ (u>0)
\]

% =====================================================
% ================= COLONNE DROITE ====================
% =====================================================

\switchcolumn
\setcounter{section}{1} %

\section{Dérivée seconde}


\subsection{Définition}

\defi Si $f$ est dérivable sur $I$, on note $f'$ sa dérivée.

Si $f'$ est dérivable sur $I$, on note $f''$ sa dérivée :  
$f''$ est appelée \textbf{dérivée seconde} de $f$.


\section{Convexité}

\subsection{Fonctions convexes}

\defi Une fonction est dite \textbf{convexe} sur un intervalle $I$ si sa courbe représentative est située \textbf{au-dessus de ses tangentes} sur $I$.

\subsection{Fonctions concaves}

\defi Une fonction est dite \textbf{concave} sur $I$ si sa courbe est située \textbf{en-dessous de ses tangentes}.

\subsection{Critère avec la dérivée seconde}

\prop Si $f$ est deux fois dérivable sur $I$ :

\[
\begin{cases}
f''(x)\ge 0 \Rightarrow f \text{ est convexe sur } I \\
f''(x)\le 0 \Rightarrow f \text{ est concave sur } I
\end{cases}
\]

\prop Équivalences importantes :

\[
f \text{ convexe } \Longleftrightarrow f' \text{ croissante}
\]
\[
f \text{ concave } \Longleftrightarrow f' \text{ décroissante}
\]



% =====================================================
% ================= PAGE SUIVANTE =====================
% =====================================================

\switchcolumn*

\subsection{Point d'inflexion}

\defi Un \textbf{point d'inflexion} est un point où la convexité de f change.

\prop Si $f''$ change de signe en $x=a$, alors la courbe admet un point d'inflexion d'abscisse $a$.

\GraphInflexion

\section{À retenir}

\begin{itemize}
\item $(f\circ u)' = u'. f'(u)$.
\item Formules usuelles : $(e^u)'$, $(\ln u)'$, $(u^n)'$, $(1/u)'$, $(\sqrt{u})'$.
\item $f''$ permet d'étudier la convexité.
\item $f''\ge 0$ : convexe ; $f''<0$ : concave.
\item Changement de signe de $f''\le 0$ $\Rightarrow$ point d'inflexion.
\end{itemize}

\end{paracol}

\end{document}
